% TOP_PROBLEM: {"id": 30007207, "title": "The numerical spacetime enclosing-area inequality in every dimension", "category_id": 16, "category": {"id": 16, "name": "physics", "display_name": "Mathematical Physics", "description": "Problems at the intersection of mathematics and physics.", "slug": "physics", "order_index": 16, "created_at": "2026-07-31T15:26:25.671Z"}, "status": "open"}
% FIRST_POSTED: 2026-10-06
% ATTEMPT_STATUS: solved
% OPENAI_CLAIM: {"schema_version":1,"source_repo":"https://github.com/openai/math","source_commit":"adc7f1241b42e322a6451854ab7e4b4c146bf78a","release_date":"2026-10-06","match":"full","resolution":"proved","independently_reviewed":false,"second_pass":"agreed","families":[{"family":"260","title":"Spacetime Penrose inequalities: enclosing area, charge, rotation, and anti-de Sitter extensions","subject":"Mathematical physics","manuscripts":[{"title":"Boundary graph deformations for the spacetime Penrose inequality","pdf_path":"preprints/Boundary-graph-deformations-for-the-spacetime-Penrose-inequality-September-27-2026/paper.pdf","readme_path":"preprints/Boundary-graph-deformations-for-the-spacetime-Penrose-inequality-September-27-2026/README.md","date":"2026-09-27","theorem_ref":"preprints/Boundary-graph-deformations-for-the-spacetime-Penrose-inequality-September-27-2026/build/neutral/00-setting.tex:124 (n3:prepared:output)"},{"title":"The spacetime Penrose inequality and enclosing area","pdf_path":"preprints/The-spacetime-Penrose-inequality-and-enclosing-area-September-27-2026/paper.pdf","readme_path":"preprints/The-spacetime-Penrose-inequality-and-enclosing-area-September-27-2026/README.md","date":"2026-09-27","theorem_ref":"preprints/The-spacetime-Penrose-inequality-and-enclosing-area-September-27-2026/build/sections/introduction.tex:87 (thm:main)"}],"lean":{"doc_path":"lean/docs/260.md","comparators":[],"declarations":[],"covers_main_theorem":false},"reasoning_trace":null}]}
\documentclass[11pt]{article}
\usepackage[margin=25mm]{geometry}
\usepackage{fontspec}
\setmainfont{Latin Modern Roman}
\usepackage{amsmath,amssymb,mathtools}
\usepackage{unicode-math}
\setmathfont{Latin Modern Math}
\usepackage{enumitem}
\usepackage{xurl,hyperref,bookmark}
\hypersetup{colorlinks=true,urlcolor=blue}
\setcounter{secnumdepth}{0}
\setlength{\parindent}{0pt}
\setlength{\parskip}{5pt}
\setlength{\emergencystretch}{3em}
\begin{document}
\section*{The numerical spacetime enclosing-area inequality in every dimension}

\subsection{OpenAI's claim}
The linked manuscripts present the following claims by OpenAI: In dimensions three and four, trapped exteriors are claimed deformable to nonnegative-scalar-curvature metrics with protected enclosing area and arbitrarily small ADM-energy loss; The proposed neutral spacetime Penrose inequality compares invariant ADM mass with minimum enclosing area in all dimensions $n\geq3$, under the listed energy, decay, and weak future-trapping hypotheses. Together they give a claimed affirmative answer in the matching scope.

\subsection{Outline of the argument}
Collars and weak boundary supports set the graph barrier. Elliptic estimates exclude contact with the floor and solve the graph-conformal deformation; its negative inner mean curvature supplies the boundary route to the neutral mass-area bound.

End preparation removes ADM momentum before a coupled graph and conformal scalar construction. Filled-equation estimates and height separation protect enclosing cuts, while the scalar identity bounds the change in mass; the Riemannian comparison then returns the inequality to the original data.

\subsection{Scope relative to this problem}
The source gives boundary deformations/prepared metric comparisons in spatial dimensions three and four. It does not itself state the target all-dimensional invariant-mass numerical inequality. The proposed neutral spacetime Penrose inequality compares invariant ADM mass with minimum enclosing area in all dimensions $n\geq3$, under the listed energy, decay, and weak future-trapping hypotheses. The question retains these same parameters and conventions. Within these assumptions, this is a claimed affirmative answer to the stated question.

\subsection{Formal verification}
For family 260, the only selected family entry is the CKS end-replacement construction and Schwarzschild equality examples. A connected-boundary inequality is conditional on the flat Penrose input; the general Bondi or spacetime inequalities are not formalized there. No Comparator is assigned to the selected manuscript here; the family page selects other papers or supporting statements. This site has not executed the Lean code. The scope descriptions do not certify the remaining manuscript arguments.

\subsection{Status caveat}
OpenAI attributes these manuscripts to work by an internal model. Its release cautions that arguments without formalization can contain mistakes. This catalogue has not independently checked the proofs; the labels report the mathematical claims attributed to OpenAI.

\subsection{Sources}
\begin{itemize}
\item OpenAI, \emph{Boundary graph deformations for the spacetime Penrose inequality}, 2026-09-27; citation key \texttt{OAI:Boundary-graph-deformations-for-the-spacetime-Penrose-inequality-September-27-2026}. \href{https://github.com/openai/math/blob/adc7f1241b42e322a6451854ab7e4b4c146bf78a/preprints/Boundary-graph-deformations-for-the-spacetime-Penrose-inequality-September-27-2026/paper.pdf}{PDF (GitHub)}; \href{https://github.com/openai/math/raw/adc7f1241b42e322a6451854ab7e4b4c146bf78a/preprints/Boundary-graph-deformations-for-the-spacetime-Penrose-inequality-September-27-2026/paper.pdf}{PDF (direct)}; \href{https://github.com/openai/math/blob/adc7f1241b42e322a6451854ab7e4b4c146bf78a/preprints/Boundary-graph-deformations-for-the-spacetime-Penrose-inequality-September-27-2026/README.md}{Manuscript page and citation}; \href{https://github.com/openai/math/blob/main/preprints/Boundary-graph-deformations-for-the-spacetime-Penrose-inequality-September-27-2026/paper.pdf}{Latest version}.
\item OpenAI, \emph{The spacetime Penrose inequality and enclosing area}, 2026-09-27; citation key \texttt{OAI:The-spacetime-Penrose-inequality-and-enclosing-area-September-27-2026}. \href{https://github.com/openai/math/blob/adc7f1241b42e322a6451854ab7e4b4c146bf78a/preprints/The-spacetime-Penrose-inequality-and-enclosing-area-September-27-2026/paper.pdf}{PDF (GitHub)}; \href{https://github.com/openai/math/raw/adc7f1241b42e322a6451854ab7e4b4c146bf78a/preprints/The-spacetime-Penrose-inequality-and-enclosing-area-September-27-2026/paper.pdf}{PDF (direct)}; \href{https://github.com/openai/math/blob/adc7f1241b42e322a6451854ab7e4b4c146bf78a/preprints/The-spacetime-Penrose-inequality-and-enclosing-area-September-27-2026/README.md}{Manuscript page and citation}; \href{https://github.com/openai/math/blob/main/preprints/The-spacetime-Penrose-inequality-and-enclosing-area-September-27-2026/paper.pdf}{Latest version}.
\item Family 260: \href{https://github.com/openai/math/blob/adc7f1241b42e322a6451854ab7e4b4c146bf78a/lean/docs/260.md}{Lean scope}.
\item \href{https://github.com/openai/math/blob/adc7f1241b42e322a6451854ab7e4b4c146bf78a/CONTENTS.md}{OpenAI catalogue entry}.
\end{itemize}
\end{document}
